\section{未来展望}

\subsection{代码的消亡与重生}

Boris 和 Dario Amodei 一年前都预言 ``到年底人们将不再亲手写代码''。到 2026 年 1 月录制节目时，Boris 表示在过去两个月中，他的所有代码（每月 200--300 个 PR）100\% 由 Claude Code 编写，他\textbf{没有手写过一行代码}。

\begin{knowledgebox}{指数思维 vs 线性思维}
Boris 强调，人类大脑天生按线性方式思考，但 AI 能力的增长是指数级的。一年前如果按线性推算，绝不可能想到模型已经能完全替代人类编码。唯一能做出正确预测的方法是``相信指数曲线，然后沿着它画下去''。
\end{knowledgebox}

\subsection{Co-Work 的平台潜力}

Greg 和 Boris 都将当前的 Agent 生态比作 iPhone App Store 的早期：

\begin{itemize}
    \item 2008 年 App Store 刚推出时，最火的是啤酒模拟器等玩具应用
    \item 没人预见到 Uber、DoorDash、TikTok 会从中诞生
    \item Co-Work 正处于类似阶段——使用场景的爆发将超出所有人的想象
\end{itemize}

Boris 预见的趋势是：越来越多的日常琐碎工作（在 App A 和 App B 之间搬运数据、格式转换、状态汇报等）将被 Agent 接管，人类将聚焦于自己真正享受的创造性工作。

\subsection{本章小结}

AI Agent 的发展遵循指数曲线。Co-Work 作为 Agent 平台的意义类似于 App Store 之于移动互联网——真正的杀手级应用尚未出现，但基础设施已经就绪。

